
Since the first order bounce-averaged drift kinetic equation may be
expressed in the same conservative form as for zero order Fokker-Planck
bounce-averaged equation, most of the boundary conditions apply in a natural
way without additional specifications, in particular at $p=0$ and $\left|
\xi _{0}\right| =1$, as a consequence of the two grids technique, one for
the fluxes and the other for the distribution function.

In matrix form, the general expression of the first order bounce averaged
drift kinetic equation 
\begin{equation}
\left\{ \mathcal{C}\left( g\right) \right\} +\left\{ \mathcal{Q}\left(
g\right) \right\} +\left\{ \mathcal{E}\left( g\right) \right\} =-\left\{ 
\mathcal{C}\left( \widetilde{f}\right) \right\} -\left\{ \mathcal{Q}\left( 
\widetilde{f}\right) \right\} -\left\{ \mathcal{E}\left( \widetilde{f}%
\right) \right\}
\end{equation}
becomes 
\begin{eqnarray}
&&\sum_{i^{\prime }=i-1}^{i^{\prime }=i+1}\sum_{j^{\prime \prime
}=j-1}^{j^{\prime \prime }=j+1}\overline{\overline{M}}_{p,l+1/2,i^{\prime
}+1/2,j^{\prime \prime }+1/2}^{\left( 0\right) }g_{0,l+1/2,i^{\prime
}+1/2,j^{\prime \prime }+1/2}^{\left( 0\right) (k)}  \nonumber \\
&&+p_{i+1/2}^{2}C^{\left( 0\right) }\left( f_{M,l+1/2,i+1/2,j+1/2},\frac{3}{2%
}\xi _{0,j+1/2}g_{l+1/2,i+1/2,j+1/2}^{\left( 0\right) (m=1)(k)}\right) 
\nonumber \\
&=&-\sum_{i^{\prime }=i-1}^{i^{\prime }=i+1}\sum_{j^{\prime \prime
}=j-1}^{j^{\prime \prime }=j+1}\widetilde{\overline{M}}_{p,l+1/2,i^{\prime
}+1/2,j^{\prime \prime }+1/2}^{\left( 0\right) }\widetilde{f}%
_{l+1/2,i^{\prime }+1/2,j^{\prime \prime }+1/2}^{\left( 0\right) }  \nonumber
\\
&&-\sum_{i^{\prime }=i-1}^{i^{\prime }=i+1}\sum_{l^{\prime }=l-1}^{l^{\prime
}=l+1}\widetilde{\overline{H}}_{\psi ,l^{\prime }+1/2,i^{\prime
}+1/2,j+1/2}^{\left( 0\right) }f_{0,l^{\prime }+1/2,i^{\prime
}+1/2,j+1/2}^{\left( 0\right) }  \nonumber \\
&&-p_{i+1/2}^{2}\widetilde{C}^{\left( 0\right) }\left(
f_{M,l+1/2,i+1/2,j+1/2},\frac{3}{2}\xi _{0,j+1/2}\widetilde{f}
_{l+1/2,i+1/2,j+1/2}^{\left( 0\right) (m=1)}\right)
\end{eqnarray}
or 
\begin{eqnarray}
&&\sum_{i^{\prime }=i-1}^{i^{\prime }=i+1}\sum_{j^{\prime \prime
}=j-1}^{j^{\prime \prime }=j+1}\overline{\overline{M}}_{p,l+1/2,i^{\prime
}+1/2,j^{\prime \prime }+1/2}^{\left( 0\right) }g_{0,l+1/2,i^{\prime
}+1/2,j^{\prime \prime }+1/2}^{\left( 0\right) (k)}  \nonumber \\
&&+p_{i+1/2}^{2}C^{\left( 0\right) }\left( f_{M,l+1/2,i+1/2,j+1/2},\frac{3}{2%
}\xi _{0,j+1/2}g_{l+1/2,i+1/2,j+1/2}^{\left( 0\right) (m=1)(k)}\right) 
\nonumber \\
&=&\widetilde{W}_{l+1/2,i+1/2,j+1/2}^{\left( 0\right) }
\end{eqnarray}
where 
\begin{eqnarray}
\widetilde{W}_{l+1/2,i+1/2,j+1/2}^{\left( 0\right) } &=&-\sum_{i^{\prime
}=i-1}^{i^{\prime }=i+1}\sum_{j^{\prime \prime }=j-1}^{j^{\prime \prime
}=j+1}\widetilde{\overline{M}}_{p,l+1/2,i^{\prime }+1/2,j^{\prime \prime
}+1/2}^{\left( 0\right) }\widetilde{f}_{l+1/2,i^{\prime }+1/2,j^{\prime
\prime }+1/2}^{\left( 0\right) (k)}  \nonumber \\
&&-\sum_{i^{\prime }=i-1}^{i^{\prime }=i+1}\sum_{l^{\prime }=l-1}^{l^{\prime
}=l+1}\widetilde{\overline{H}}_{\psi ,l^{\prime }+1/2,i^{\prime
}+1/2,j+1/2}^{\left( 0\right) }f_{0,l^{\prime }+1/2,i^{\prime
}+1/2,j+1/2}^{\left( 0\right) (k)}  \nonumber \\
&&-p_{i+1/2}^{2}\widetilde{C}^{\left( 0\right) }\left(
f_{M,l+1/2,i+1/2,j+1/2},\frac{3}{2}\xi _{0,j+1/2}\widetilde{f}
_{l+1/2,i+1/2,j+1/2}^{\left( 0\right) (m=1)}\right)
\end{eqnarray}

Here, $\left( k\right) $ is the iteration index number, while $C^{\left(
0\right) }$ and $\widetilde{C}^{\left( 0\right) }$ are the electron-electron
self-collision operators that make the equation non-linear.

\subsubsection{Internal boundaries for function g$^{\left( 0\right) }$}

In the weak collision ``banana'' regime, it is shown in Sec.\ref
{lowcollision} that $g^{\left( 0\right) }\left( \psi ,p,\xi _{0}\right) =0$
for trapped electrons. Therefore, the drift kinetic equation must be solved
in a reduced domain of integration, as for $f_{0}^{\left( 0\right) }\left(
\psi ,p,\xi _{0}\right) ,$ but with different and more simple internal
boundary conditions. Indeed, the only constraint is to enforce the condition 
$g^{\left( 0\right) }\left( \psi ,p,\xi _{0}\right) =0$ for $\left| \xi
_{0}\right| \leq \xi _{0T}.$

%TCIMACRO{
%\TeXButton{Fig_implicitdomaing}{\fig{figures/Implicitdomaing.eps}{Trapped domain for the first order distribution g}{Implicitdomaing}{10cm}}}
%BeginExpansion
\fig{figures/Implicitdomaing.eps}{Trapped domain for the first order distribution g}{Implicitdomaing}{10cm}%
%EndExpansion

Let define grid points 
\begin{equation}
\left\{ 
\begin{array}{l}
1:\left( l+1/2,i+1/2,j_{T,l+1/2}^{+}-1/2\right)  \\ 
2:\left( l+1/2,i+1/2,j_{T,l+1/2}^{+}+1/2\right)  \\ 
3:\left( l+1/2,i+1/2,j_{T,l+1/2}^{-}+1/2\right)  \\ 
4:\left( l+1/2,i+1/2,j_{T,l+1/2}^{-}-1/2\right) 
\end{array}
\right. 
\end{equation}
as indicated in Fig.\ref{Implicitdomaing}. Here $j_{T,l+1/2}^{-}$ is the
index that corresponds to the trapped/passing boundary $\xi _{0}=-\xi _{0T},$
while $j_{T,l+1/2}^{+}$ corresponds to the symmetric boundary with respect
to the axis $\xi _{0}=0,$ i.e. $\xi _{0}=\xi _{0T}.$

Since $g^{\left( 0\right) }$ is prescribed in the trapped domain, all the
region may be removed from the calculations, and the reduction of the domain
of integration in the momentum phase space leads to introduce a new matrix $%
\widehat{\widehat{M}}_{p}^{\left( 0\right) }$ whose coefficients are defined
as follows 
\begin{equation}
\widehat{\widehat{M}}_{p,l+1/2,i+1/2,j^{\prime }+1/2}^{\left( 0\right) }=%
\overline{\overline{M}}_{p,l+1/2,i+1/2,j+1/2}^{\left( 0\right) }
\end{equation}
with $j^{\prime }=j$, when $0\leq j<j_{T,l+1/2}^{-},$ and $j^{\prime
}=j-m_{\xi _{0T,l+1/2}},$ when $j_{T,l+1/2}^{+}<j\leq n_{\xi _{0}}$. Here, $%
m_{\xi _{0T,l+1/2}}=j_{T,l+1/2}^{+}-j_{T,l+1/2}^{-}$ corresponds to the
number of grid points removed in the pitch-angle direction, and with this
definition, it is easy to cross-check that $j_{T,l+1/2}^{\prime
+}-j_{T,l+1/2}^{\prime -}=j_{T,l+1/2}^{+}-j_{T,l+1/2}^{-}-m_{\xi
_{0T,l+1/2}}=0.$ Since grid points $2$ and $4$ are not connected, the
following condition applies 
\begin{equation}
\left\{ 
\begin{array}{l}
\widehat{\widehat{M}}_{p,l+1/2,i-1/2,\left[ j_{T,l+1/2}^{\prime }-1/2\right]
+1}^{\left( 0\right) }=0 \\ 
\widehat{\widehat{M}}_{p,l+1/2,i+1/2,\left[ j_{T,l+1/2}^{\prime }-1/2\right]
+1}^{\left( 0\right) }=0 \\ 
\widehat{\widehat{M}}_{p,l+1/2,i+3/2,\left[ j_{T,l+1/2}^{\prime }-1/2\right]
+1}^{\left( 0\right) }=0
\end{array}
\right.
\end{equation}
and conversely 
\begin{equation}
\left\{ 
\begin{array}{l}
\widehat{\widehat{M}}_{p,l+1/2,i-1/2,\left[ j_{T,l+1/2}^{\prime }+1/2\right]
-1}^{\left( 0\right) }=0 \\ 
\widehat{\widehat{M}}_{p,l+1/2,i+1/2,\left[ j_{T,l+1/2}^{\prime }+1/2\right]
-1}^{\left( 0\right) }=0 \\ 
\widehat{\widehat{M}}_{p,l+1/2,i+3/2,\left[ j_{T,l+1/2}^{\prime }+1/2\right]
-1}^{\left( 0\right) }=0
\end{array}
\right.
\end{equation}
where $j_{T,l+1/2}^{\prime -}=j_{T,l+1/2}^{\prime +}=j_{T,l+1/2}^{\prime }.$
As internal boundaries match already existing diagonals of the matrix $%
\overline{\overline{M}}_{p,l+1/2,i+1/2,j+1/2}^{\left( 0\right) },$ $\widehat{%
\widehat{M}}_{p,l+1/2,i+1/2,j^{\prime }+1/2}^{\left( 0\right) }$ remains a
simple nine diagonal matrix, and is therefore less complex than the $15$
diagonals matrix $\widehat{\overline{M}}_{p,l+1/2,i+1/2,j^{\prime
}+1/2}^{\left( 0\right) }$ corresponding to the Fokker-Planck equation.

\subsubsection{Internal boundaries for $\widetilde{f}$}

Since no inversion is required for calculating the vector $\widetilde{W}%
_{l+1/2,i+1/2,j+1/2}^{\left( 0\right) },$ it is consequently evaluated at
all grid points, including those corresponding to trapped electron domain.
Consequently, the reduction along the pitch-angle direction that results
from the condition $g^{\left( 0\right) }\left( \psi ,p,\xi _{0}\right) =0$
is only performed after the calculation of $\widetilde{W}%
_{l+1/2,i+1/2,j+1/2} $ over the full $\xi _{0}$ grid. Once done, the
prescription on the index number $j^{\prime }$ that holds for the matrix
conversion 
\begin{equation}
\overline{\overline{M}}_{p,l+1/2,i+1/2,j+1/2}^{\left( 0\right) }\rightarrow 
\widehat{\widehat{M}}_{l+1/2,i+1/2,j+1/2p,l+1/2,i+1/2,j^{\prime
}+1/2}^{\left( 0\right) }
\end{equation}
may be then applied in a similar way for $\widetilde{W}_{l+1/2,i+1/2,j+1/2},$
which leads to define a new vector $\widehat{\widetilde{W}}%
_{p,l+1/2,i+1/2,j^{\prime }+1/2}^{\left( 0\right) }$ whose size is reduced 
\begin{equation}
\widetilde{W}_{l+1/2,i+1/2,j+1/2}^{\left( 0\right) }\rightarrow \widehat{%
\widetilde{W}}_{p,l+1/2,i+1/2,j^{\prime }+1/2}^{\left( 0\right) }
\end{equation}
and finally the drift kinetic equation on the reduced grid is simply given
by the relation 
\begin{equation}
\sum_{i^{\prime }=i-1}^{i^{\prime }=i+1}\sum_{j^{\prime \prime }=j^{\prime
}-1}^{j^{\prime \prime }=j^{\prime }+1}\widehat{\widehat{M}}%
_{p,l+1/2,i+1/2,j^{\prime \prime }+1/2}^{\left( 0\right)
}g_{0,l+1/2,i^{\prime }+1/2,j^{\prime \prime }+1/2}^{\left( 0\right) (k)}=%
\widehat{\widetilde{W}}_{p,l+1/2,i+1/2,j^{\prime }+1/2}^{\left( 0\right) }
\end{equation}
where all the trapped electron domain is removed.

\subsubsection{External boundaries}

The problem of external boundaries discussed in Sec. \ref{FPboundaries} for
the zero-order Fokker-Planck equation may be applied in a straightforward
manner to the first-order drift kinetic ones, since it may be expressed in a
similar conservative form. Consequently, if the zero-order Fokker-Planck
equation conserves the total number of particles without specific
conditions, the solution $g$ may be found without additional constraints at
the boundaries of the integration domain. Conversely, if Dirichlet
conditions are considered for the Fokker-Planck equation, like enforcing the
Maxwellian solution at $p=0$, similar and consistent conditions must be also
applied for the drift kinetic equation.
